\documentclass{article}

% if you need to pass options to natbib, use, e.g.:
%     \PassOptionsToPackage{numbers, compress}{natbib}
% before loading neurips_2026

% The authors should use one of these tracks.
% Before accepting by the NeurIPS conference, select one of the options below.
% 0. "default" for submission
\usepackage{xcolor} 
\usepackage{neurips_2026}
% the "default" option is equal to the "main" option, which is used for the Main Track with double-blind reviewing.
% 1. "main" option is used for the Main Track
%  \usepackage[main]{neurips_2026}
% 2. "position" option is used for the Position Paper Track
%  \usepackage[position]{neurips_2026}
% 3. "eandd" option is used for the Evaluations & Datasets Track
% \usepackage[eandd]{neurips_2026}
% if you need to opt-in for a single-blind submission in the E&D track:
%\usepackage[eandd, nonanonymous]{neurips_2026}
% 4. "creativeai" option is used for the Creative AI Track
%  \usepackage[creativeai]{neurips_2026}
% 5. "sglblindworkshop" option is used for the Workshop with single-blind reviewing
% \usepackage[sglblindworkshop]{neurips_2026}
% 6. "dblblindworkshop" option is used for the Workshop with double-blind reviewing
%  \usepackage[dblblindworkshop]{neurips_2026}

% After being accepted, the authors should add "final" behind the track to compile a camera-ready version.
% 1. Main Track
% \usepackage[main, final]{neurips_2026}
% 2. Position Paper Track
%  \usepackage[position, final]{neurips_2026}
% 3. Evaluations & Datasets Track
% \usepackage[eandd, final]{neurips_2026}
% 4. Creative AI Track
%  \usepackage[creativeai, final]{neurips_2026}
% 5. Workshop with single-blind reviewing
%  \usepackage[sglblindworkshop, final]{neurips_2026}
% 6. Workshop with double-blind reviewing
%  \usepackage[dblblindworkshop, final]{neurips_2026}
% Note. For the workshop paper template, both \title{} and \workshoptitle{} are required, with the former indicating the paper title shown in the title and the latter indicating the workshop title displayed in the footnote.
% For workshops (5., 6.), the authors should add the name of the workshop, "\workshoptitle" command is used to set the workshop title.
% \workshoptitle{WORKSHOP TITLE}

% "preprint" option is used for arXiv or other preprint submissions
 % \usepackage[preprint]{neurips_2026}


\usepackage[utf8]{inputenc} % allow utf-8 input
\usepackage[T1]{fontenc}    % use 8-bit T1 fonts
\usepackage{hyperref}       % hyperlinks
\usepackage{url}            % simple URL typesetting
\usepackage{booktabs}       % professional-quality tables
\usepackage{amsfonts}       % blackboard math symbols
\usepackage{nicefrac}       % compact symbols for 1/2, etc.
\usepackage{microtype}      % microtypography
        % colors
\usepackage{amsmath}
\usepackage{amssymb}
\usepackage{mathtools}
\usepackage{amsthm}
\usepackage{thmtools}
\usepackage[capitalize,noabbrev]{cleveref}
\usepackage{xargs}
\usepackage{natbib} 
\usepackage{datatool}
\usepackage{tikz}
\usepackage{xparse}
\DeclareMathOperator*{\argmax}{arg\,max}
\DeclareMathOperator*{\argmin}{arg\,min}

\newcommand{\E}{\mathbb{E}}
\newcommand{\prob}{\mathbb{P}}

\newcommand{\Var}{\mathrm{Var}}
\newcommand{\R}{\mathbb{R}}
\newcommand{\spaceg}{G^*}
\newcommand{\simplexg}{\ensuremath{\bar{G^*}}}
\newcommand{\localrad}[2]{\ensuremath{R_\ndata(#1, #2)}}
\def\state{\overrightarrow{X}} % this means that the main state variable will always be X. We use capital X to the random variable 
\def\bara{\bar{\alpha}}
\def\ausiliar{\rho}
\def\barg{\bar{g}}
\def\barsigma{\bar{\sigma}}
\def\lstate{\operatorname{x}} % this means that the main state variable will always be x. We use lower x to indicate values that the random variable takes, as for densities for example.
\def\rmd{\mathrm{d}} %
\def\law{\mathcal{L}}
\def\interval{[0,T]}
\def\dt{\rmd \operatorname{t}}
\def\dbrown{\rmd\!\brown}
\def\brown{\operatorname{B}}
\def\rset{\mathbb{R}}
\def\rsetpos{\mathbb{R}_{\geq 0}}
\def\statedim{d}
\def\targetdist{\mu_{*}}
\def\eqsp{\;}
\def\score{\nabla \log \rho}
\def\realint{\int_{\R^d}}


%Useful commands from my papers with Sylvain Eric and others
\newcommand{\kl}[2]{\operatorname{D}_{\operatorname{KL}}\left(#1||#2\right)}
\newcommand{\W}[2]{\operatorname{W}\left(#1||#2\right)}
\newcommand{\dist}[2]{\operatorname{D}\left(#1||#2\right)}
\newcommand{\dens}[2]{#1\left(#2\right)}
\newcommand{\condd}[3]{\dens{#1}{#3 | #2}}
\newcommand{\utility}[1]{\mathrm{U}(#1)}
\newcommand{\mseloss}[2]{\operatorname{MSE}(#1; #2)}
\newcommand{\measop}[1]{f(#1)}
\def\measnoise{\varepsilon_\lmeas}

\newcommand{\fwmarg}[1]{p_{#1}} % distribution of the foward sde at instant #1
\newcommand{\fwmargd}[2]{\dens{p_{#1}}{#2}} % densitiy of the foward sde at instant #1
\newcommand{\fwtrans}[2]{p_{#2 | #1}} %transition kernel for foward SDE

\newcommand{\mudata}{\mu_{\text{data}}}


\newcommand{\PE}[2]{\mathbb{E}_{#1}\left[#2\right]} % Expected value operator
\newcommand{\PP}[1]{\mathbb{P}\left(#1\right)} % Expected value operator
\newcommand{\PV}[2]{\mathbb{V}_{#1}\left[#2\right]} % Variance operator
\newcommand{\CPE}[3]{\PE{#1}{#3 | #2}} % Conditional expected value

\newcommand{\gausspdf}{\varphi}
\newcommand{\rme}{\mathrm{e}}
\newcommand{\ola}{\overleftarrow}
\newcommand{\ora}{\overrightarrow}
\newcommand{\eqdef}{:=}
\newcommand{\param}{\theta}
\newcommand{\paramdim}{d_\theta}
\def\paramsp{\Theta}
\newcommand{\pihat}{\widehat{\pi}_N^{(\beta,\param)}}
\newcommand{\sched}{\beta}
\def \Nc{\mathcal{N}}
\newcommand{\pdata}{\pi_{data}}


\newcommand{\logp}[1]{\log \tilde{p}_{T - #1} (\x{#1})}
\newcommand{\pforward}[1]{\vec{
p}_{#1}}
\newcommand{\backtildex}[1]{\tilde{X}_{#1}}

\newcommand{\trainedsolution}[1]{X^{\param}_{#1}}
\newcommand{\trainedsolutioninit}[2]{X^{\theta,#2}_{#1}}
\newcommand{\trainedscorex}[1]{s_{\param}(T- #1, \back{#1})}
\newcommand{\trainedscorenov}[1]{s_{\param}(T- #1, \back{#1}^x)}
\newcommand{\dl}{\, d\lambda_{\mathbb{R}^d}}
\newcommand{\ptildex}{\tilde{p}_t(x)}
\newcommand{\pforwardx}[1]{\vec{p}_{#1}}

\newcommand{\gam}{\gamma(x)}
\newcommand{\trainedscore}[1]{s_{\param}(T- #1, \trainedsolution{#1})}

\newcommand{\back}[1]{\overleftarrow{X}_{#1}}
\newcommand{\backscoretilde}[1]{\nabla \log \tilde{p}_{T - #1} (\back{#1})}

\newcommand{\backtildedensityW}[1]{\tilde{p}^{\param}_{#1}}
\newcommand{\backscore}[1]{\nabla \log \pforward{T - #1} (\back{#1})}
\newcommand{\backscorenov}[1]{\nabla \log \pforward{T - #1} (\back{#1}^x)}

\newcommand{\backscoretheta}[1]{\nabla \log \pforward{T - #1} (\trainedsolution{#1})}

\newcommand{\scorep}[1]{\nabla \log \pforward{#1}}
\newcommand{\net}[1]{s_{\param}(T-#1)}

\newcommand{\discsolution}[2][\param]{p_{#2}^{#1}}
\newcommand{\schedulerVE}{g}
\newcommand{\information}{\mathcal{I}}
\newcommand{\pinf}{\pi_{\infty}}

\newcommand{\ind}{\mathds{1}}
\newcommand{\Id}{\mathrm{I}}
\newcommand{\scfc}{\mathsf{s}}

\newcommand{\initerror}{\mathsf{E}_{\operatorname{init}}}
\newcommand{\trainerror}{\mathsf{E}_{\operatorname{train}}}
\newcommand{\discerror}{\mathsf{E}_{\operatorname{disc}}}
\newcommand{\rmL}{\mathrm{L}}
\newcommand{\calF}{\mathcal{F}}
\newcommand{\borelians}[1]{\mathcal{B}(#1)}
\def\tensprod{\otimes}
\def\ndata{n}
\makeatletter
\NewDocumentCommand{\distn}{g g}{%
  d_{\ndata}\IfValueTF{#1}{\left( #1 , #2 \right)}{}%
}
\def\maxray{r_{\ndata}}
\NewDocumentCommand{\distnparam}{g g}{%
  \bar{d}_{\ndata}\IfValueTF{#1}{\left( #1 , #2 \right)}{}%
}
\newcommand{\nballs}[3]{N(#1, #2, #3)}
\newcommand{\chunk}[3]{{#1}^{#2:#3}}
\newcommand{\emin}[1]{\widehat{#1}^{\ndata}}
\newcommand{\tmin}[1]{#1^{\star}}
\NewDocumentCommand{\mfun}{o o}{\operatorname{m}\IfValueTF{#1}{(#2; #1)}{}}
\NewDocumentCommand{\zeromfun}{o o}{\operatorname{g}_{\IfValueTF{#1}{#1}{}}\IfValueTF{#2}{\left(#2\right)}{}}
\NewDocumentCommand{\eintmfun}{o o}{\operatorname{M}_{\ndata}\IfValueTF{#2}{\IfValueTF{#1}{(#2; #1)}{(#2)}}{}}
\NewDocumentCommand{\tintmfun}{o}{\operatorname{M}_{\star}\IfValueTF{#1}{(#1)}{}}
\NewDocumentCommand{\empmeas}{o o}{\mu_{\operatorname{emp}}\IfValueTF{#1}{\IfValueTF{#2}{\left(#1; #2\right)}{\left(#1\right)}}{}}
\NewDocumentCommand{\lpnorm}{m o o}{\left\|#1\right\|_{\operatorname{L}^{#2}(#3)}}
\NewDocumentCommand{\onorm}{m o}{\left\|#1\right\|_{\psi_#2}}
\NewDocumentCommand{\flow}{o o}{\operatorname{T}\IfValueTF{#1}{_{#1}}{}\IfValueTF{#2}{(#2)}{}}
\NewDocumentCommand{\jflow}{o o}{\operatorname{J}_{\operatorname{T}\IfValueTF{#1}{_{#1}}{}}\IfValueTF{#2}{(#2)}{}}
\NewDocumentCommand{\invflow}{o o}{\operatorname{T}^{-1}\IfValueTF{#1}{_{#1}}{}\IfValueTF{#2}{(#2)}{}}
\NewDocumentCommand{\pset}{m}{\mathcal{P}\left(#1\right)}
\NewDocumentCommand{\mlip}{o}{\operatorname{b}\IfValueTF{#1}{\left(#1\right)}{}}
\newcommand{\logdet}[1]{\operatorname{logdet}\left(#1\right)}
\newcommand{\diameter}[1]{\operatorname{Diam}(#1)}
\def\refnfmeas{u}
\newcommand{\pushforward}[2]{{#1}_{\# #2}}
\newcommand{\worstconcentrationvar}[2]{\operatorname{Z}_{\ndata}(#1, #2)}
\newcommand{\gabriel}[1]{\textcolor{red}{GC: #1}}
\newcommand{\tiz}[1]{\textcolor{red}{TF: #1}}
\newcommand{\cvx}[1]{\overline{#1}}
\newcommand{\lrmc}[1]{\cvx{R}_{\ndata}(#1)}
\newcommand{\bernoulli}[1]{\operatorname{Bern}(#1)}
\def\ctesupconcentration{\overline{\operatorname{C}}}
\newcounter{hypH}
\newenvironment{hypH}{
    \refstepcounter{hypH}
    \begin{itemize}
    \item[{\bf H\arabic{hypH}}]
    }
{\end{itemize}}

% Note. For the workshop paper template, both \title{} and \workshoptitle{} are required, with the former indicating the paper title shown in the title and the latter indicating the workshop title displayed in the footnote. 
\title{Generalizing Score based Generative Model Applications}

\author{%
Tiziano, Sylvain, Gabriel, Thomas\thanks{Use footnote for providing further information
about author (webpage, alternative address)---\emph{not} for acknowledging
funding agencies.} \\
Department of Computer Science\\
Cranberry-Lemon University\\
Pittsburgh, PA 15213 \\
\texttt{hippo@cs.cranberry-lemon.edu} \\
% examples of more authors
% \And
% Coauthor \\
% Affiliation \\
% Address \\
% \texttt{email} \\
% \AND
% Coauthor \\
% Affiliation \\
% Address \\
% \texttt{email} \\
% \And
% Coauthor \\
% Affiliation \\
% Address \\
% \texttt{email} \\
% \And
% Coauthor \\
% Affiliation \\
% Address \\
% \texttt{email} \\
}


\begin{document}
\theoremstyle{plain}
\newtheorem{theorem}{Theorem}[section]
\newtheorem{proposition}[theorem]{Proposition}
\newtheorem{lemma}[theorem]{Lemma}
\newtheorem{corollary}[theorem]{Corollary}
\theoremstyle{definition}
\newtheorem{definition}[theorem]{Definition}
\newtheorem{assumption}[theorem]{Assumption}
\theoremstyle{remark}
\newtheorem{remark}[theorem]{Remark}
\crefname{assumption}{assumption}{assumptions}
\maketitle
\begin{abstract}
It is not clear in the literature if SDE based diffusion models can be used for any distribution. 
In particular it is not clear if the theoretical backward equation is well posed in any case and if
the standard approximation using a trained score and a gaussian initialization is 
sufficient for approximating any distribution. 
In recent literature in particular a debate around the possibility of applying diffusion models to heavytail distributions is grown. 
In this paper we theoretically show that the use of early stopping and the replacement of the gaussian initialization with a suitable one
allows the generalization of the framework of diffusion models to any distribution. In particular we show that the backward is well posed and we generalize a KL-based theorem 
to any existent distribution by replacing the gaussian initialiation, that is not sufficient, with a suitable one. 
We test our idea on toy datasets, classical datasets and heavytail datasets, showing that this framework is efficient in the number of steps and improves the generation in the heavytail context. 
\end{abstract}
\section{Introduction}
In recent years, generative models have emerged as a cornerstone of machine learning due to their capacity to sample high-dimensional distributions.
Score-based generative models (SGMs) were introduced in their discrete form by Ho et al.~\citep{ho2020denoising} and in their continuous formulation by Song et al.~\citep{song2021scorebased}, building on several years of pioneering work~\citep{sohl2015deep, song2019generative}.
Since then, they have achieved state-of-the-art performance in image and video generation, surpassing GANs, VAEs, and Normalizing Flows \citep{dhariwal2021diffusionbeatGan, muller2022medfusion}. Their frameworks form the foundation of Stable Diffusion \citep{rombach2022high} and many of the most prominent publicly available models today, including DALL·E 3 \citep{Betker2023dalle3}, Imagen 3 \citep{Baldridge2024Imagen3}, and recent advanced architectures like FLUX \citep{blackForestLabs_FLUX1dev_2025}. SGMs have also demonstrated competitive results in audio \citep{chen2020wavegrad, kong2021diffwave} and text generation \citep{li2022diffusionlm, zou2023survey}, positioning them as a flexible and powerful framework across diverse data modalities.

The core idea of SGMs is twofold: the data distribution is gradually corrupted by adding Gaussian noise, and a neural network is trained to reverse this process, recovering the original sample from its corrupted version. New samples are then generated by starting from standard Gaussian noise and iteratively applying the trained network to progressively denoise it.

Theoretical work on SGMs pursues a twofold goal: understanding \textit{how} the approximate backward diffusion converges to the distribution of interest, and proving \textit{that} it does. While it is clear that approximating the score function and using an approximate distribution both introduce errors in generation, it remains an open question under which conditions these approximations are sufficient to guarantee convergence to the true distribution. To establish such guarantees, existing works typically rely on strong assumptions on the data distribution. For instance, \citep{debortoli2022convergence} restricts to bounded support, \citep{gao2025wasserstein} requires log-concavity, and \citep{conforti2023kl} assumes a finite moment of order eight. More broadly, all existing theoretical analyses of SGMs impose some form of strong assumption on the data distribution \citep{bruno2024wasserstein, silveri2025beyond, bortoli2021diffusion, strasman2025analysis}.

A more fundamental issue concerns the backward SDE itself. The founding paper on continuous SGMs \citep{song2021scorebased} relies on \citep{anderson1982reverse} to introduce it, yet \citep{anderson1982reverse} establishes this result only for distributions with bounded support. While other works \citep{bortoli2021diffusion} invoke alternative references \citep{haussmann1986time}, it remains unclear what minimal assumptions on the data distribution are required to rigorously justify the backward SDE, and therefore the diffusion framework itself.

Understanding for which distributions the diffusion framework can be rigorously introduced, and when convergence can be established, is thus a question of both theoretical and practical importance. On the practical side, the recent surge in high-quality image generation has prompted researchers to explore whether diffusion can be extended to heavy-tailed distributions. Several approaches have been proposed, including normalizing flows \citep{hickling2025flexibletails, mcdonald2022cometflows, laszkiewicz2022marginal}, variational autoencoders \citep{lafon2023vae}, GANs \citep{allouche2022evgan, huster2021pareto}, and diffusion-based methods \citep{pandey2025heavytailed, yoon2023levy, shariatian2025heavy}, yet none has fully solved the problem. Several works \citep{pandey2025heavytailed, shariatian2025heavy} explicitly highlight the difficulties SGMs face in this regime, but it remains unclear whether these failures stem from fundamental limitations of the framework or from practical issues such as training instability, poor initialization, or SDE discretization.

In this paper, we show that early stopping and proper initialization together allow the diffusion framework to be applied, theoretically, to any distribution. The idea of improving generation through non-standard initialization has already been explored \citep{zheng2023truncated, zand2024diffusion, DBLP:journals/corr/abs-2306-12360, lee2022priorgrad, shen2025information, scholz2026warm}, but always with the goal of improving efficiency or generation quality. Our contribution is fundamentally different: we provide a theoretical guarantee that diffusion models can be applied to \textit{any} distribution in the sense of KL convergence, without any assumption on the latter.

Precisely, this paper makes the following contributions.
\begin{itemize}
    \item We prove that the backward SDE is well-posed for any data distribution, establishing that the diffusion framework is theoretically grounded in full generality.
    \item We prove that combining early stopping with a flexible initialization scheme that replaces the standard Gaussian prior is sufficient to guarantee convergence of the approximate diffusion to the target distribution without any assumption on the latter, generalizing \cite[Theorem 1]{conforti2023kl} to this end.
    \item We validate our initialization-aware approach on standard benchmarks, demonstrating that a well-chosen initialization improves generation quality while reducing the number of required denoising steps.
    \item We take a first step toward modeling heavy-tailed distributions with standard diffusion models. We discuss what families of distributions are suitable choices for the initialization prior, and provide empirical evidence on toy datasets, synthetic climatic data, and real climatic datasets that proper initialization meaningfully improves over standard diffusion in this regime.
\end{itemize}
%\section{Related Works}
The idea of improving generation through a non-standard initialization has already been explored in the literature, though from a different perspective than ours. \citep{zheng2023truncated} proposed a diffusion-based generative model that shortens the diffusion trajectory by learning, via an Adversarial Autoencoder (AAE), a prior at a truncated diffusion time that approximates the aggregated posterior of the forward noising process. \citep{zand2024diffusion} introduced DiNof, which uses a Glow Normalizing Flow to approximate the noised distribution at an arbitrary diffusion step, thereby reducing sampling time. \citep{DBLP:journals/corr/abs-2306-12360} proposed Discrete Walk-Jump sampling, combining a score-based model and an Energy-Based Model with a single noise level to stabilize training and accelerate sampling. Further works along this line include \citep{lee2022priorgrad, shen2025information, scholz2026warm}.
All these works share a common motivation: improving generation efficiency or quality. Our contribution is fundamentally different in nature. We show that proper initialization, combined with early stopping, provides a theoretical guarantee that diffusion models can be applied to \textit{any} distribution, in the sense of KL convergence, without any assumption on the latter.
\paragraph{SGM Guarantees. }
To the best of our knowledge, existing theoretical works on SGMs establish convergence and error guarantees exclusively under a fixed Gaussian initialization and under . Such guarantees are developed either within the KL-divergence framework \citep{conforti2023kl} or within the Wasserstein-distance literature \citep{debortoli2022convergence,strasman2025analysis,strasman2025wasserstein,gao2025wasserstein,bruno2024wasserstein,silveri2025beyond}, and uniformly rely on this standard initialization assumption, without considering alternative initialization schemes. Our theoretical analysis focuses on the initialization error, decoupled from the training of the score model, and introduces a theoretically grounded criterion for optimizing the initialization of the backward process.
\paragraph{Complementary Methods. } 
Flow matching methods \citep{lipman2023flow} have recently shifted focus to transport cost and toward designing informative priors \citep{issachar2024designing, tong2024improving, kornilov2024optimal, feng2025on}. For example, Issachar et al. \citep{issachar2024designing} constructs conditional priors from training-data centroids in order to improve generation. Tong et al \citep{tong2024improving} proposes a batch algorithm for optimal transport flow matching.
Our approach provides a general prior-design framework that supports both unconditional and conditional priors. 
By using priors closer to the target distribution, it offers the potential to shorten the transport path.
Alternatively, one-step methods like Consistency Models \citep{song2023consistency} and MeanFlow \citep{geng2025mean} reduce sampling time, but learning a direct map from a Gaussian prior to complex data remains challenging. Using intermediate distributions offers a promising approach for more efficient one-step trajectories.
\section{Generalized Variance Exploding diffusion}
We discuss in this section the existence of the forward and backward diffusion equations in the context of Variance Exploding framework. 
Let $(\Omega, \calF, \prob)$ be a probability space and $\rmL^2(\Omega, \calF, \prob)$ the space of square-integrable random variables on $\Omega$. We denote by $\mudata$ the unknown data distribution on $\R^d$, and the goal of SGMs is to generate new samples from $\mudata$ given only an observed dataset.
\subsection{Forward and Backward Process}
Given $\state_0 \sim \mudata$ the distribution of interest and the standard forward Variance Exploding (VE) equation \citep{song2021scorebased}
\begin{equation}\label{eq:VE_SDE}
\rmd \state_t = \schedulerVE_t \, \dbrown_t\eqsp.
\end{equation}
where $\schedulerVE_t : \R^+ \to \R$ is assumed to be bounded and infinitely differentiable, we can state by applying \citet[Theorem 2.5]{karatzas1991brownian} without any assumption on the data distribution that the forward solution $(\state_t)_{t\in[0,T]}$ is defined on all intervall for $T>0$.
Defining 
\[
\sigma_t^2 \coloneqq \int_0^t g_s \, \rmd s\eqsp.
\]
It can be stated that the solution
\[
\state_t = \state_0 + \int_{0}^{t} \schedulerVE_s ^2 \rmd s \overset{\mathcal{L}}{=} \state_0 + \sigma_t Z\eqsp,
\]
Therefore as well known in the literature the solution simply results to be the convolution of the data distribution with a gaussian. 
In particular, for $t>0$, $\state_t$ admits a density $\pforward{t}(x)$ with respect to the Lebesgue measure $\lambda_{\R^d}$ and this density is highly regular. In particular the function $ (t,x)\mapsto\pforward{t}(x) \in C^{\infty}([\delta, T] \times \R^d)$, as shown in Lemma ---- . 

Concerning the backward equation the situation, might seem more complicated. The existence of the backward equation is not given for any distribution and any forward sde admitting an explicit solution. In fact the sufficient condition stated in reference articles \citep{anderson1982reverse,haussmann1986time} require the data distribution 
or the density of the forward (when available) to respect some regularity conditions. Thankfully in the (VE) context where the forward solution is simply the convolution of the distributin with a gaussian, this doesn't represent a problem as stated in our first proposition. 
\begin{proposition}
Given $\mudata$ a distribution and $\state_t, \pforward{t}(x)$ the solutions of the \eqref{eq:VE_SDE} the backward equation 
\begin{equation}\label{eq:VE_back_sde}
\rmd \back{t} =   \barg_t^2 \backscore{t} \rmd t + \barg_t \, \rmd \tilde{\brown}_t \eqsp ,
\end{equation}
where $\barg_t \coloneqq \schedulerVE_{T-t}$ and $\tilde{\brown}_t$ is an independent Brownian motion.
For all initialization $\back{0}$ and for all $\delta>0$  \eqref{eq:VE_back_sde} admits a strong and unique solution $\back{} =\{\back{t}\}_{[0,T-\delta]}$ on the interval $[0, T-\delta]$. 
In particular if $\back{0} \sim \pforward{T}$, for all $\delta>0$ it holds $\back{} \overset{\mathcal{L}}{=}  \{\state\}_{[\delta, T]}$.
\end{proposition}
The proposition that we presented establishes the existence of a backward strong solution for all

\subsection{KL control on approximated process}
In practice,  $\mu_*$ is not available; instead, we only have access to an i.i.d. sample $\state_0^i \sim \mu_*$, $1 \leq i \leq n$. 
Consequently, the solution of \cref{eq:VE_back_sde} that would enable sample generation is not directly accessible, since neither the score nor the terminal distribution $\pforwardx{T}$ of the backward process are known. 
Furthermore, the backward SDE does not admit a closed-form solution, making it necessary to rely on discretization schemes.

For the purpose of theoretical analysis, following common practice, we may rely on either the Euler–Maruyama scheme \cite{gao2025wasserstein, silveri2025beyond} or the Exponential Euler–Maruyama scheme \cite{conforti2023kl, strasman2025analysis}. Remarkably, in the VE setting these two schemes coincide, allowing for a unified treatment of the KL bound.

\paragraph{Score approximation and discretization. }
By replacing the real score function $\nabla \log \pforward{T - t}$ using a trained neural network $s_{\param}(T- t,\cdot)$, we
introduce the discretized backward dynamics
\begin{equation}\label{eq:disceq}
\rmd \trainedsolution{t}=  \barg_t^2 \trainedscore{t_k} \, \rmd t + \barg_t \, \rmd\tilde{\brown}_t\eqsp,
\end{equation}
for $t \in [t_k, t_{k+1}]$, where $t_{k+1} = t_k + h_{k+1}$, with $t_0 = 0$ and $t_N = T-\delta$
which can be equivalently written, for  $t \in [0, T - \delta]$, as
\begin{equation}\label{eq:disceq}
\rmd \trainedsolution{t}
    =  \sum_{k=0}^{N-1}\barg_t^2 \trainedscore{t_k} 1_{t \in [t_k, t_{k+1}]} \, \rmd t 
        + \barg_t \, \rmd\tilde{\brown}_t\eqsp.
\end{equation}
This discretization scheme results in the following stochastic sampling dynamics:
\begin{equation}\label{eq:discsolution}
\trainedsolution{t_{k+1}}= \trainedsolution{t_{k}} + (\sigma^2_{T-t_{k}}-\sigma^2_{T-t_{k+1}}) \trainedscore{t_k} + \sqrt{\sigma^2_{T-t_{k}}-\sigma^2_{T-t_{k+1}}}Z \eqsp,
\end{equation}
for $k = 0,...,N-1$, $Z \sim \mathcal{N}(0, \Id_d)$.
The only assumption we have to make to establish our result is a score approximation assumption : 
\begin{itemize}
    \item \(\discsolution{0}\) and $\lambda_{\R^d}$ are mutually absolutely continuous.
    \item $\kl{\pforward{t}}{\discsolution[\param]{0}} <  \infty$
\end{itemize}
\begin{hypH}
\label{assum:train}
There exists $\varepsilon>0$ such that for all $0 \leq k \leq N-1$,
$$
\E \Big[\|\backscore{t_k} - s_{\param}(T-t_k, \back{t_k})\|^2\Big] \leq \varepsilon\eqsp.
$$
\end{hypH}
We say that 
$$\int_{\mathbb{R}^d} \log\frac{p_T}{p_\param^0}(x ) p_T(x) dx$$
is finite only in certain scenarios. Therefore : gaussian is not sufficient, the family of $p_\param$ has at least to respect : 
$$\sup_{\param \in \param} \int_{\mathbb{R}^d} \log\frac{p_T}{p_\param^0}(x ) p_T(x) dx < + \infty$$

Passando per l informazione di Fisher riesco a dimostrare che é sempre nulla. 
\section{Experiments}
\begin{itemize}
    \item We present experiments with normal datasets. We show that everything works well and we also saved some time. This is not new so this part will not be central. 
    \item We present experiments with heavytail distribution. We show that it works pretty well, so that there is hope and ongoing work to do to correctly choose the flow $p_\theta$ which is exactly the point of research where we are stuck. 
    \item We can include Nicolas and he could help me with doing the simulations on the heavytail datasets. As datasets i propose the synthetic dataset i created and the Comephere data which can be cleaned by me and Nicolas. 
\end{itemize}

\bibliographystyle{plainnat} 
\bibliography{biblio}
\end{document}